%----------------------------------------------------------------------
% Reiniciar numeração das figuras que aparecem no apêndice
\setcounter{figure}{0}

\chapter{Código-fonte das implementações}
\label{apend:codigo}

\vspace{-1.9cm}

O código-fonte completo, os arquivos de resultados em formato CSV e os scripts de
geração dos gráficos estão disponíveis no repositório do trabalho. Reproduzem-se a
seguir os trechos essenciais das implementações.

\section{Insertion Sort}

\begin{lstlisting}[language=C++,caption={Ordenação por inserção do subvetor \texttt{v[esq..dir]}.}]
void insertionSort(std::vector<int>& v, int esq, int dir, Contadores& c) {
    for (int i = esq + 1; i <= dir; ++i) {
        const int chave = v[i];
        int j = i - 1;
        // Desloca para a direita todos os elementos maiores que 'chave'.
        while (j >= esq) {
            ++c.comparacoes;
            if (v[j] <= chave) break;
            v[j + 1] = v[j];
            ++c.trocas;
            --j;
        }
        if (j + 1 != i) {
            v[j + 1] = chave;
        }
    }
}
\end{lstlisting}

\section{Partição de Hoare}

\begin{lstlisting}[language=C++,caption={Partição em torno do valor do pivô. Ao final, \texttt{v[esq..j]} contém os elementos menores ou iguais ao pivô e \texttt{v[i..dir]}, os maiores ou iguais.}]
static void particao(std::vector<int>& v, int esq, int dir,
                     int& i, int& j, int pivo, Contadores& c) {
    i = esq;
    j = dir;
    do {
        while (true) { ++c.comparacoes; if (v[i] >= pivo) break; ++i; }
        while (true) { ++c.comparacoes; if (v[j] <= pivo) break; --j; }
        if (i <= j) {
            std::swap(v[i], v[j]);
            ++c.trocas;
            ++i;
            --j;
        }
    } while (i <= j);
}
\end{lstlisting}

\section{Escolha do pivô}

\begin{lstlisting}[language=C++,caption={Mediana de três valores, com contagem de comparações.}]
static int mediana3(int a, int b, int d, Contadores& c) {
    ++c.comparacoes;
    if (a > b) std::swap(a, b);   // garante a <= b
    ++c.comparacoes;
    if (b > d) {                  // a <= b e d < b: mediana e max(a,d)
        ++c.comparacoes;
        return (a > d) ? a : d;
    }
    return b;                     // a <= b <= d
}

static int escolherPivo(const std::vector<int>& v, int esq, int dir,
                        EstrategiaPivo p, Contadores& c) {
    const int meio = esq + (dir - esq) / 2;   // evita overflow de (esq+dir)/2
    switch (p) {
        case EstrategiaPivo::PRIMEIRO: return v[esq];
        case EstrategiaPivo::MEIO:     return v[meio];
        case EstrategiaPivo::MEDIANA3: return mediana3(v[esq], v[meio], v[dir], c);
    }
    return v[meio];
}
\end{lstlisting}

\section{Motor único das três versões}

\begin{lstlisting}[language=C++,caption={Motor parametrizado pelo corte \texttt{M} e pela estratégia de pivô.}]
static void ordena(std::vector<int>& v, int esq, int dir,
                   Contadores& c, int M, EstrategiaPivo p) {
    if (esq >= dir) return;

    // Interrompe a particao para subvetores com MENOS de M elementos.
    // Com M <= 1 nenhum subvetor nao trivial e interrompido, e o algoritmo
    // degenera na versao recursiva pura.
    if (dir - esq + 1 < M) {
        insertionSort(v, esq, dir, c);
        return;
    }

    int i, j;
    const int pivo = escolherPivo(v, esq, dir, p, c);
    particao(v, esq, dir, i, j, pivo, c);
    ordena(v, esq, j, c, M, p);
    ordena(v, i, dir, c, M, p);
}

// As tres versoes exigidas sao atalhos com parametros fixos:
void quicksortRecursivo(std::vector<int>& v, Contadores& c, EstrategiaPivo p) {
    quicksort(v, c, 1, p);
}
void quicksortHibrido(std::vector<int>& v, Contadores& c, int M) {
    quicksort(v, c, M, EstrategiaPivo::MEIO);
}
void quicksortHibridoMediana3(std::vector<int>& v, Contadores& c, int M) {
    quicksort(v, c, M, EstrategiaPivo::MEDIANA3);
}
\end{lstlisting}
